\documentclass[12pt]{article}
\usepackage[a4paper,margin=1in]{geometry}
\usepackage{graphicx}
\usepackage[export]{adjustbox}
\usepackage{hyperref}
\usepackage{fontspec}
\setmainfont{DejaVu Sans}
\usepackage{amsmath}
\usepackage{amssymb}
\usepackage{listings}

\begin{document}
VIỆN CÔNG NGHỆ APEX — KHOA MÁY TÍNH
Thạc sĩ Khoa học Dữ liệu & Trí tuệ Nhân tạo (MS-DSAI)
Năm học: 2026–2027
• Trình độ đào tạo: Sau đại học (Thạc sĩ Khoa học)
• Tổng số tín chỉ: 36 Tín chỉ Mỹ / 90 ECTS

1. Tóm tắt Tổng quan & Tầm nhìn Chương trình
Sứ mệnh Chương trình: Chương trình Thạc sĩ Khoa học Dữ liệu & Trí tuệ Nhân tạo (MS-DSAI) tại Viện Công nghệ Apex là một chương trình đào tạo chuyên nghiệp nâng cao, được thiết kế nhằm khỏa lấp khoảng cách giữa nền tảng máy tính lý thuyết và triển khai công nghiệp mang lại tác động cao. Chương trình giảng dạy trang bị cho các học viên tốt nghiệp nền tảng toán học nghiêm ngặt, kỹ thuật hệ thống máy học có khả năng mở rộng và năng lực quản trị AI đạo đức cần thiết để dẫn dắt các cuộc chuyển đổi phức tạp dựa trên dữ liệu trong các doanh nghiệp toàn cầu hiện đại.

36
TỔNG SỐ TÍN CHỈ

2 Năm
THỜI GIAN TOÀN THỜI GIAN

98%
TỶ LỆ CÓ VIỆC LÀM

1:8
TỶ LỆ GIẢNG VIÊN/HỌC VIÊN

2. Mục tiêu Giáo dục của Chương trình (PEOs)
Sau khi hoàn thành xuất sắc chương trình MS-DSAI, các học viên tốt nghiệp dự kiến sẽ đạt được các cột mốc sự nghiệp và chuyên môn sau trong vòng 3 đến 5 năm sau khi tốt nghiệp:

PEO 1 (Thành thạo Kỹ thuật): Dẫn dắt việc kiến trúc, thiết kế và triển khai các hệ thống trí tuệ nhân tạo phức tạp và các nền tảng dữ liệu doanh nghiệp xử lý các đường ống tính toán đa terabyte.

PEO 2 (Đổi mới Sáng tạo Ứng dụng): Đề xuất các giải pháp thuật toán nguyên bản cho các vấn đề thực tế, mơ hồ trong các lĩnh vực như y tế, tài chính, hệ thống tự hành và logistics hiện đại.

PEO 3 (Lãnh đạo Đạo đức): Thiết lập các khung quản trị AI vững chắc, tính minh bạch của thuật toán, bảo mật quyền riêng tư và tuân thủ đạo đức trong các tổ chức công nghệ.

PEO 4 (Hợp tác Xuất sắc): Truyền đạt rõ ràng các góc nhìn định lượng và tính toán phức tạp tới các bên liên quan cấp quản lý không thuộc lĩnh vực kỹ thuật, các đội ngũ kỹ thuật liên chức năng và các cơ quan quản lý.

3. Các Trụ cột Chính của Chương trình Giảng dạy

1. Toán học Nghiêm ngặt
Bao quát sâu về giải tích vectơ, xác suất nhiều chiều, tối ưu hóa lồi, các quá trình ngẫu nhiên và đại số tuyến tính làm nền tảng cho các kiến trúc học sâu hiện đại.

2. Kỹ thuật Hệ thống
Kinh nghiệm thực hành với tính toán phân tán, dòng dữ liệu thời gian thực, cụm tăng tốc GPU, MLOps, đóng gói container và microservices.

3. AI Có trách nhiệm
Tích hợp trực tiếp sự công bằng của thuật toán, khả năng giải thích (XAI), học liên kết, luật bảo mật dữ liệu và sự tuân thủ an toàn theo quy định trong tất cả các chuyên ngành.

Thạc sĩ Khoa học Dữ liệu & Trí tuệ Nhân tạo — Sổ tay Học thuật Trang 1 / 12

4. Cấu trúc Chương trình Giảng dạy & Phân bổ Tín chỉ
Bằng MS-DSAI yêu cầu hoàn thành 36 tín chỉ sau đại học. Học viên tiến triển qua một trình tự có cấu trúc bao gồm Nền tảng Cốt lõi (18 tín chỉ), Môn tự chọn Chuyên ngành (9 tín chỉ), Thực tập Công nghiệp / Chuyên đề Sau đại học (3 tín chỉ), và Luận văn Thạc sĩ hoặc Đề án Tốt nghiệp (6 tín chỉ).

Phân loại
Số tín chỉ yêu cầu
Tương đương ECTS
Mô tả

Nền tảng Cốt lõi
18 Tín chỉ
45 ECTS
6 môn học cốt lõi bắt buộc bao gồm toán học, thuật toán, học máy, học sâu, hệ thống phân tán, và NLP/CV.

Môn tự chọn Chuyên ngành
9 Tín chỉ
22.5 ECTS
3 môn học chuyên sâu được chọn từ một chuyên ngành chỉ định (Hệ thống ML, Kiến trúc Dữ liệu Đám mây, hoặc Quản trị AI).

Chuyên đề Sau đại học / Thực tập
3 Tín chỉ
7.5 ECTS
Chuỗi chương trình hướng dẫn từ doanh nghiệp, phương pháp nghiên cứu và các hội thảo chủ đề mới nổi.

Luận văn / Đề án Tốt nghiệp
6 Tín chỉ
15 ECTS

Luận văn nghiên cứu gốc hai học kỳ hoặc đồ án tốt nghiệp doanh nghiệp.
Tổng Yêu cầu Bằng cấp
36 Tín chỉ
90 ECTS
4 Học kỳ (24 Tháng Toàn thời gian / 36 Tháng Bán thời gian)
5. Lộ trình Tiến trình Học thuật
Sơ đồ dưới đây thể hiện trình tự học tập được đề xuất cho sinh viên cao học toàn thời gian trong suốt bốn học kỳ:
HỌC KỲ 1 (10 TC)
HỌC KỲ 2 (10 TC)
HỌC KỲ 3 (9 TC)
HỌC KỲ 4 (7 TC)
DSAI-501 (3 TC)
Toán ứng dụng & Thống kê
DSAI-502 (3 TC)
Cấu trúc dữ liệu & Python
DSAI-503 (3 TC)
Cơ sở Học máy
DSAI-500 (1 TC) Hội thảo khoa học
DSAI-504 (3 TC)
Học sâu & Kiến trúc
DSAI-505 (3 TC)
Hệ thống Dữ liệu lớn Phân tán
DSAI-506 (3 TC)
Thị giác máy tính & NLP
DSAI-590 (1 TC) Phòng thí nghiệm Đạo đức
Chuyên đề Tự chọn I (3 TC)
Chuyên sâu Chuyên ngành
Chuyên đề Tự chọn II (3 TC)
Chuyên sâu Chuyên ngành
DSAI-698 (3 TC)
Đồ án tốt nghiệp / Luận văn Giai đoạn I
Chuyên đề Tự chọn III (3 TC)
Chuyên sâu Chuyên ngành
DSAI-699 (3 TC)
Luận văn / Đồ án tốt nghiệp II
Bảo vệ Tốt nghiệp & Thi vấn đáp
DSAI-600 (1 TC) Hồ sơ năng lực
6. Tóm tắt các Yêu cầu Tốt nghiệp Chính
Điểm TBM Tích lũy Tối thiểu: Sinh viên phải duy trì Điểm Trung bình Môn (GPA) tích lũy tối thiểu là 3.00 trên 4.00 cho tất cả các môn học cao học đã đăng ký.
Tiêu chuẩn Năng lực Cốt lõi: Yêu cầu đạt điểm 'B' (3.0) trở lên trong tất cả sáu môn học Cơ sở Cốt lõi.
Yêu cầu Học tại cơ sở: Ít nhất 27 tín chỉ phải được hoàn thành trực tiếp hoặc đồng bộ chính thức tại Viện Công nghệ Apex (Apex Institute of Technology).
Bảo vệ Luận văn / Đồ án tốt nghiệp: Bảo vệ công khai thành công và được hội đồng phê duyệt Luận văn Thạc sĩ hoặc Đồ án tốt nghiệp Doanh nghiệp.
Thạc sĩ Khoa học Dữ liệu & Trí tuệ Nhân tạo — Sổ tay Học thuật
Trang 2 / 12

7. Đề cương các Môn học Cốt lõi — Năm 1, Học kỳ 1
Học kỳ đầu tiên thiết lập các nguyên lý toán học nền tảng và năng lực tính toán hiệu năng cao cần thiết cho nghiên cứu và ứng dụng học máy nâng cao.
DSAI-501 — Xác suất Ứng dụng & Thống kê Toán học cho Khoa học Dữ liệu
3 Tín chỉ • Bắt buộc
Tổng quan Môn học: Trình bày nâng cao về lý thuyết phân phối xác suất đa biến, ước lượng cực đại giống, suy luận Bayesian, Chuỗi Markov, thuật toán tối đa hóa kỳ vọng (EM), kiểm định giả thuyết và không gian vectơ đa chiều.
Kết quả Học tập: Xây dựng các mô hình xác suất chặt chẽ; triển khai các phương pháp lấy mẫu Monte Carlo; suy ra các hàm mất mát dạng đóng để ước lượng tham số.
Lý thuyết Xác suất Suy luận Bayesian Chuỗi Markov Giải tích Vectơ
Đánh giá: Bài tập về nhà (25\%), Bài thi Giữa kỳ (35\%), Bài thi Lý thuyết Cuối kỳ (40\%).
DSAI-502 — Cấu trúc Dữ liệu Thuật toán Nâng cao & Hệ thống Python
3 Tín chỉ • Bắt buộc
Tổng quan Môn học: Lập trình cấp hệ thống chuyên sâu bằng Python và C++ cho kỹ thuật dữ liệu. Bao gồm phân tích bộ nhớ tiệm cận, tính toán mảng vectơ hóa, xử lý đồ thị, xử lý song song và các mô hình thiết kế phần mềm hướng đối tượng.
Kết quả Học tập: Xây dựng các cấu trúc dữ liệu thuật toán tối ưu bộ nhớ; tối ưu hóa mức sử dụng bộ nhớ cho các phép biến đổi ma trận; viết các tiện ích mở rộng C song song cho môi trường Python.
Python C-API Cơ chế Nội bộ NumPy Thuật toán Đồ thị Phân tích Bộ nhớ
Đánh giá: Bài thực hành Lập trình định kỳ hai tuần một lần (40\%), Đồ án Kiến trúc Hệ thống (30\%), Bài thi Cuối kỳ (30\%).
DSAI-503 — Cơ sở Học máy & Học thống kê
3 Tín chỉ • Bắt buộc
Tổng quan Môn học: Phân tích toán học toàn diện về các thuật toán học có giám sát và học không giám sát.

Các chủ đề bao gồm hồi quy tuyến tính, hồi quy logistic, SVM, tập hợp cây quyết định (Random Forests, XGBoost), giảm chiều dữ liệu (PCA, t-SNE), và các kỹ thuật chuẩn hóa (L1/L2, ElasticNet).
Kết quả học tập: Cài đặt các thuật toán học máy cốt lõi từ các nguyên lý toán học mà không phụ thuộc vào các thư viện cấp cao; đánh giá sự đánh đổi giữa độ chệch và phương sai (bias-variance trade-offs) bằng phương pháp toán học.
Học có giám sát Tối ưu hóa lồi Phương pháp tập hợp Ranh giới thống kê
Đánh giá: Bài tập lớn xây dựng thuật toán (35\%), Bài kiểm tra giữa kỳ (30\%), Bài tập lớn \& Báo cáo cuối kỳ (35\%).
Thạc sĩ Khoa học Dữ liệu \& Trí tuệ Nhân tạo — Sổ tay Học thuật
Trang 3 / 12

8. Đề cương các môn học cốt lõi — Năm 1, Học kỳ 2
Học kỳ thứ hai nâng cao vào các kiến trúc trí tuệ nhân tạo hiện đại, hạ tầng dữ liệu lớn phân tán và xử lý thông tin đa thức phức tạp.
DSAI-504 — Kiến trúc Học sâu \& Mạng Nơ-ron
3 Tín chỉ • Cốt lõi
Tổng quan môn học: Nghiên cứu sâu về mạng nơ-ron nhân tạo, các công cụ đạo hàm tự động, mạng tích chập (CNN), các mô hình chuỗi (RNN, LSTM), và cơ chế Transformer tự chú ý (self-attention). Phát triển mô hình thực hành sử dụng PyTorch.
Kết quả học tập: Thiết kế các kiến trúc nơ-ron tùy chỉnh; cài đặt lan truyền ngược từ đầu (from scratch); tối ưu hóa chiến lược tìm kiếm siêu tham số trên các cụm GPU phân tán.
PyTorch Transformers Cơ chế chú ý Giảm độ dốc
Đánh giá: Phòng thí nghiệm đánh giá hiệu năng học sâu (30\%), Thử thách kiến trúc mô hình (30\%), Bài thi cuối kỳ (40\%).
DSAI-505 — Hệ thống Dữ liệu lớn Phân tán \& Điện toán Đám mây
3 Tín chỉ • Cốt lõi
Tổng quan môn học: Kỹ thuật xây dựng các đường ống phân tích có thể mở rộng trên các cụm đám mây và biên. Nghiên cứu sâu về Apache Spark, Hadoop MapReduce, hệ thống tập tin phân tán (HDFS/S3), cơ sở dữ liệu SQL phân tán và các công cụ xử lý luồng (Kafka, Flink).
Kết quả học tập: Kiến trúc các tác vụ xử lý dữ liệu phân tán có khả năng chịu lỗi; tối ưu hóa các đồ thị thực thi DAG; truy vấn các tập dữ liệu hàng terabyte dưới ràng buộc độ trễ dưới một giây.
Apache Spark Kafka Streaming Kiến trúc Đám mây SQL Phân tán
Đánh giá: Bài thực hành kỹ thuật cụm phân tán (40\%), Nghiên cứu tình huống về khả năng mở rộng (25\%), Bài tập lớn cuối kỳ (35\%).
DSAI-506 — Thị giác Máy tính Áp dụng \& Xử lý Ngôn ngữ Tự nhiên
3 Tín chỉ • Cốt lõi
Tổng quan môn học: Tập trung kép vào trí tuệ thị giác và ngôn ngữ. Bao gồm phát hiện đối tượng, phân đoạn ngữ nghĩa, các Mô hình Ngôn ngữ Lớn (LLM) hiện đại, phân đoạn từ (tokenization), không gian nhúng (embedding spaces), tạo ra tăng cường truy xuất (RAG), và các phương pháp tinh chỉnh (fine-tuning).
Kết quả học tập: Tinh chỉnh các mô hình thị giác-ngôn ngữ đã được huấn luyện trước; triển khai các đường ống RAG trên các kho tài liệu doanh nghiệp; đánh giá các chỉ số căn chỉnh ngữ nghĩa.
LLMs \& RAG Phát hiện Đối tượng Cơ sở dữ liệu Vector Hugging Face
Đánh giá: Hackathon AI đa thức (35\%), Tái tạo bài báo nghiên cứu (25\%), Bài thi cuối kỳ (40\%).
Thạc sĩ Khoa học Dữ liệu \& Trí tuệ Nhân tạo — Sổ tay Học thuật
Trang 4 / 12

9. Chuyên ngành 1: Kỹ thuật Hệ thống Học máy
Trọng tâm chuyên ngành: Được thiết kế cho các sinh viên có nguyện vọng trở thành Kỹ sư MLOps chính, Kiến trúc sư Hệ thống AI và Trưởng nhóm Nền tảng Học máy. Tập trung vào vòng đời sản xuất, tăng tốc phần cứng, triển khai mô hình có thể mở rộng, đánh giá liên tục và giám sát hạ tầng.
DSAI-610 — MLOps Có thể Mở rộng \& Hạ tầng Sản xuất
3 Tín chỉ • Tự chọn Chuyên ngành
Tích hợp liên tục / Triển khai liên tục (CI/CD) cho Học máy, kho tính năng (Feast), mô hình

các kho lưu trữ mô hình (MLflow), phục vụ mô hình trên Kubernetes, tự động mở rộng quy mô động, phát hiện độ lệch mô hình (model drift) và các chiến lược tự động khôi phục (rollback).
DSAI-611 — Học tăng cường & Các tác tử tự hành
3 Tín chỉ • Môn tự chọn Hướng chuyên sâu
Quy trình quyết định Markov (MDP), Q-learning, các phương pháp Policy Gradient, Deep Q-Network (DQN), Proximal Policy Optimization (PPO), học tăng cường đa tác tử và RLHF (Học tăng cường từ phản hồi của con người) để tinh chỉnh căn chỉnh.
DSAI-612 — AI Tạo sinh & Kỹ thuật Mô hình Ngôn ngữ Lớn
3 Tín chỉ • Môn tự chọn Hướng chuyên sâu
Các đổi mới kiến trúc trong Mạng đối kháng tạo sinh (GAN), Mô hình khuếch tán (Diffusion Models), các biến thể Transformer, tinh chỉnh hiệu quả tham số (LoRA, QLoRA), lượng thể hóa (AWQ, GGUF) và các công cụ suy luận hiệu năng cao (vLLM, TensorRT-LLM).
Chồng Hạ tầng & Phần cứng Tính toán của Hướng chuyên sâu 1
Sinh viên theo học Hướng chuyên sâu 1 nhận được các token tính toán cấp phát riêng trên Cụm Siêu máy tính GPU Apex AI (8 nút NVIDIA H100 mỗi pod), giúp trải nghiệm thực hành huấn luyện phân tán đa nút.
Chồng Điều phối & Phục vụ
Kubernetes & Kubeflow Pipelines
Triton Inference Server
Quản lý Cụm Ray
Tối ưu hóa & Lưu trữ
ONNX Runtime & TensorRT
Cơ sở dữ liệu Vector Pinecone & Milvus
Weave & Weights & Biases
Thạc sĩ Khoa học Dữ liệu & Trí tuệ Nhân tạo — Sổ tay Học thuật
Trang 5 / 12

10. Hướng chuyên sâu 2: Hệ thống Dữ liệu & Phân tích Đám mây
Trọng tâm Hướng chuyên sâu: Được thiết kế cho các Kiến trúc sư Nền tảng Dữ liệu, Giám đốc Dữ liệu (CDO) và Trưởng nhóm Phân tích Đám mây tương lai. Tập trung vào các đường ống dữ liệu băng thông cao, kho dữ liệu doanh nghiệp, kiến trúc truyền dòng sự kiện thời gian thực và thiết kế data lakehouse hiện đại.
DSAI-620 — Các hệ cơ sở dữ liệu hiệu năng cao & NoSQL
3 Tín chỉ • Môn tự chọn Hướng chuyên sâu
Công cụ lưu trữ nội bộ (B-Trees, LSM-Trees), phục hồi WAL, độ cô lập giao dịch ACID, các giao thức đồng thuận phân tán (Raft, Paxos), cơ sở dữ liệu key-value, cơ sở dữ liệu tài liệu (MongoDB), cơ sở dữ liệu wide-column (Cassandra) và cơ sở dữ liệu đồ thị (Neo4j).
DSAI-621 — Các hệ thống truyền dòng thời gian thực & Kiến trúc Hướng sự kiện
3 Tín chỉ • Môn tự chọn Hướng chuyên sâu
Tính nhị nguyên dòng-bảng (Stream-table duality), các microservice hướng sự kiện, xử lý dòng có trạng thái với Apache Flink, ngữ nghĩa xử lý chính xác một lần (exactly-once), xử lý sự kiện sai thứ tự với watermark và kiến trúc sự kiện phát hiện gian lận tốc độ cao.
DSAI-622 — Kho dữ liệu, Lakehouse & BI Doanh nghiệp
3 Tín chỉ • Môn tự chọn Hướng chuyên sâu
Kiến trúc chồng dữ liệu hiện đại sử dụng Snowflake, Databricks Delta Lake và Apache Iceberg. Mô hình hóa chiều, sơ đồ Ngôi sao/Bông tuyết (Star/Snowflake schemas), các framework chuyển đổi dbt, tầng ngữ nghĩa và kiểm thử chất lượng dữ liệu tự động.
Các nghiên cứu điển hình về hệ thống công nghiệp được đúc kết
Ngành công nghiệp nền tảng
Điểm nổi bật về kiến trúc
Mục tiêu kỹ thuật cốt lõi
Dịch vụ Tài chính
Luồng dữ liệu thị trường dưới một mili giây
Truyền dòng sự kiện không mất dữ liệu với Kafka & Flink
Quy mô Thương mại Điện tử
Data Lakehouse cấp độ Petabyte
Các giao dịch ACID trên bộ lưu trữ đối tượng với Apache Iceberg
Lô-gi-stíc Toàn cầu
Công cụ đo đạc từ xa IoT thời gian thực
Phân vùng cơ sở dữ liệu chuỗi thời gian & truy vấn động
Thạc sĩ Khoa học Dữ liệu & Trí tuệ Nhân tạo — Sổ tay Học thuật
Trang 6 / 12

11. Hướng chuyên sâu 3: Quản trị AI, Đạo đức & Công nghệ Trách nhiệm
Trọng tâm Hướng chuyên sâu: Chuẩn bị cho các nhà lãnh đạo kỹ thuật, cố vấn pháp lý, kiểm toán viên AI và các chuyên gia chính sách. Tập trung vào kiểm toán độ công bằng của thuật toán, tuân thủ pháp lý (EU AI Act, NIST AI RMF), AI có thể giải thích (XAI) và học máy bảo tồn quyền riêng tư.
DSAI-630 — Chính sách AI, Khung Pháp lý & An toàn Thuật toán

3 Tín chỉ • Môn tự chọn Chuyên ngành
Phân tích sâu về pháp luật toàn cầu bao gồm Đạo luật AI của EU (EU AI Act), Khung Quản lý Rủi ro AI của NIST Hoa Kỳ, và các tiêu chuẩn bảo mật dữ liệu toàn cầu (GDPR, CCPA). Phân loại rủi ro, các giao thức kiểm toán kỹ thuật, báo cáo tuân thủ và trách nhiệm pháp lý.
DSAI-631 — Học máy Bảo vệ Quyền riêng tư & Hệ thống Liên kết
3 Tín chỉ • Môn tự chọn Chuyên ngành
Các nguyên lý toán học của Quyền riêng tư Vi phân (DP), Mã hóa Đồng hình (HE), Tính toán Dạng bên Bảo mật (SMPC), các thuật toán gom nhóm học liên kết (FedAvg, FedProx), và phòng thủ chống lại các cuộc tấn công suy luận thành viên.
DSAI-632 — AI Có thể Giải thích (XAI) & Tính Công bằng Thuật toán
3 Tín chỉ • Môn tự chọn Chuyên ngành
Cơ chế gán thuộc tính đặc trưng (SHAP, LIME, Integrated Gradients), giải thích phản thực tế, các định nghĩa toán học về tính công bằng thuật toán (demographic parity, equalized odds), và giảm thiểu độ lệch (bias) trong quy trình huấn luyện.
Chuỗi Hội thảo Nghiên cứu & Báo cáo chuyên đề từ Doanh nghiệp
Chuyên ngành 3 bao gồm các buổi diễn thuyết nghiên cứu bắt buộc định kỳ hai tuần một lần với các nhà đạo đức học hàng đầu, chuyên gia pháp lý và giám đốc tuân thủ cấp cao từ các cơ quan quản lý quốc tế và các công ty công nghệ.
Chuỗi Diễn thuyết A: AI Pháp lý
Tập trung vào việc vận hành các yêu cầu tài liệu kỹ thuật của Đạo luật AI của EU trong các triển khai rủi ro cao (y tế, chấm điểm tín dụng, tuyển dụng).
Chuỗi Diễn thuyết B: Kiểm toán Kỹ thuật
Các buổi thực hành về triển khai bộ thử nghiệm tính công bằng tự động trong quy trình tích hợp liên tục bằng cách sử dụng các công cụ mã nguồn mở (AIF360, Fairlearn).
Thạc sĩ Khoa học Dữ liệu & Trí tuệ Nhân tạo — Sổ tay Học thuật
Trang 7 / 12

12. Khung Luận văn Thạc sĩ & Đồ án Tốt nghiệp Capstone
Là yêu cầu hoàn thành chương trình bằng MS-DSAI, các học viên phải hoàn thành Hướng Luận văn Nghiên cứu Hàn lâm nguyên bản hoặc Hướng Đồ án Capstone Doanh nghiệp do doanh nghiệp tài trợ trong Học kỳ 3 và 4 (tổng cộng 6 tín chỉ).
Lựa chọn A: Luận văn Nghiên cứu Hàn lâm
Đối tượng: Sinh viên có kế hoạch học tiếp Tiến sĩ hoặc làm việc tại các phòng nghiên cứu.
Đóng góp nghiên cứu lý thuyết hoặc thuật toán nguyên bản.
Hướng dẫn bởi cố vấn giảng viên + 2 thành viên hội đồng.
Yêu cầu nộp một bài báo phản biện độc lập cho một hội nghị hàng đầu (NeurIPS, ICML, KDD, CVPR, ACL).
Lựa chọn B: Đồ án Capstone Doanh nghiệp
Đối tượng: Sinh viên hướng tới các vị trí lãnh đạo ngay trong ngành công nghiệp.
Dự án được tài trợ giải quyết thách thức thực tế của doanh nghiệp.
Đồng hướng dẫn bởi một kỹ sư chính từ doanh nghiệp & cố vấn giảng viên.
Bàn giao kho lưu trữ mã nguồn sẵn sàng đưa vào sản xuất + bảo vệ trước hội đồng lãnh đạo.
13. Tiêu chí Đánh giá & Tiến độ Cột mốc
Giai đoạn Cột mốc
Kế hoạch Dự kiến
Sản phẩm Bàn giao Chính
Tỷ trọng
Bảo vệ Đề cương
Cuối Học kỳ 3
(Tuần 14)
Tổng quan Tài liệu Chính thức, Phát biểu Vấn đề, Đề xuất Phương pháp luận & Sự phê duyệt của Hội đồng Cố vấn.
15%
Đánh giá Tiến độ Giữa kỳ
Học kỳ 4 (Tuần 7)
Chạy bản thử nghiệm thực thi (prototype), đánh giá chuẩn so sánh sơ bộ, kiểm duyệt kiến trúc mã nguồn.
25%
Báo cáo Luận văn/Đồ án
Học kỳ 4 (Tuần
13)
Bản thảo nghiên cứu đạt chất lượng xuất bản hoặc cuốn tài liệu kỹ thuật toàn diện.
35%
Bảo vệ Báo cáo Công khai
Học kỳ 4 (Tuần
15)
Thuyết trình công khai (30 phút) + phản biện chất vấn nghiêm ngặt từ hội đồng đánh giá của khoa sau đại học.
25%
Sơ đồ Trực quan Tiến độ Dự án
1
Đề cương
HK 3, T14
2
Bản Alpha
HK 4, T7
3
Nộp Bản thảo
HK 4, T13
4
Bảo vệ Báo cáo
HK 4, T15
Thạc sĩ Khoa học Dữ liệu & Trí tuệ Nhân tạo — Sổ tay Học thuật
Trang 8 / 12

14. Thang điểm Học thuật & Tiêu chuẩn Đánh giá

Apex Institute of Technology enforces strict academic evaluation criteria to maintain global accredited standards.
Graduate performance is graded according to the letter-grade system below:
Grade
Grade
Points
Percentage
Equivalent
Performance Definition
A
4.00
93% – 100%
Exemplary Mastery. Demonstrates superior theoretical understanding and
flawless practical execution.
A-
3.70
90% – 92%
Excellent Performance. Outstanding comprehension with minor non-critical
technical errors.
B+
3.30
87% – 89%
Very Good. Solid technical execution above average graduate standards.
B
3.00
83% – 86%
Good / Satisfactory. Meets standard requirements for core course progression.
B-
2.70
80% – 82%
Below Graduate Average. Credited, but triggers Academic Progress review if
cumulative GPA falls < 3.00.
C+ /
C
2.30 / 2.00
73% – 79%
Unsatisfactory Performance. Course must be repeated if listed as a mandatory
core course.
F
0.00
Below 70%
Failing Grade. Zero credit awarded. Mandatory repeat required.
15. Academic Integrity & Responsible Generative AI Policy
Strict Honor Code Policy: Academic dishonesty, including uncredited plagiarism, code fabrication, contract
cheating, and unauthorized examination assistance, results in immediate referral to the Graduate Disciplinary
Board with potential academic expulsion.
Policy on Generative AI Usage in Coursework
Permitted Use (with Citation): Utilizing LLM assistants for syntax debugging, spellchecking, and brainstorming is
permitted ONLY when explicitly authorized in writing by the course instructor. All AI-assisted code segments or text must
be explicitly disclosed in a dedicated "Generative AI Attribution Index" appended to the submission.
Prohibited Use: Generating complete assignment codebases, submitting synthetic research literature reviews, or
employing AI during closed-book theoretical examinations constitutes severe academic misconduct.
16. Academic Probation & Remediation Protocols
Any student whose cumulative GPA falls below 3.00 at the conclusion of any semester is automatically placed on
Academic Probation. The student is assigned an Academic Advisory Specialist and granted one semester to raise
their cumulative GPA to 3.00 or higher. Failure to satisfy probation terms leads to formal academic dismissal.
Master of Science in Data Science & Artificial Intelligence — Academic Handbook
Page 9 of 12

17. Distinguished Program Faculty
The MS-DSAI program is delivered by internationally recognized scholars, researchers, and former principal industry
scientists:
Prof. Elena Rostova, Ph.D.
Department Chair & Professor of AI
Ph.D. Computer Science (Stanford)
Specialization: High-dimensional optimization, deep generative
architectures, and federated learning protocols. Former
Research Lead at Google Brain.
Prof. Marcus Vance, Ph.D.
Director, Center for Autonomous Systems
Ph.D. Robotics (MIT)
Specialization: Computer vision, real-time spatial sensing, deep
reinforcement learning for robotics, and autonomous agent
navigation.
Dr. Aisha Al-Hassan, Ph.D.
Associate Professor of Data Systems
Ph.D. Computer Engineering (UC Berkeley)
Specialization: Distributed database systems, stream
processing, and scalable cloud analytics architectures. Co-
inventor of key stream indexing algorithms.
Dr. Jonathan Thorne, J.D., Ph.D.
Director, Responsible AI & Policy Lab
Ph.D. CS (Oxford), J.D. (Harvard)
Specialization: Algorithmic fairness, explainable AI, technical
auditing frameworks, and international regulatory compliance
law.

18. Các Trung tâm & Cơ sở Nghiên cứu Liên kết
Trung tâm Siêu máy tính Apex
Cụm máy tính hiện đại sở hữu 64 nút GPU
NVIDIA H100 với kết nối InfiniBand 100Gbps,
dành riêng cho nghiên cứu của sinh viên và
tiền huấn luyện mô hình quy mô lớn.
Phòng Thí nghiệm Kiểm thử AI Có trách nhiệm
Phòng thí nghiệm liên ngành tập trung vào việc
phát triển các bộ công cụ phát hiện thiên vị tự động,
bộ chuẩn đánh giá độ bền đối kháng và các
mạng học liên kết bảo vệ quyền riêng tư.
Cơ sở Thị giác Tự hành
Khu vực tính toán không gian và robot tiên tiến
được trang bị các hệ thống ghi hình chuyển động
đa máy ảnh, hệ thống LiDAR và môi trường
mô phỏng GPU biên.
Các Khoản Tài trợ Nghiên cứu Bên ngoài Đang Hoạt động
Tài trợ NSF Frontier (#2025-AI-889): "Xác minh Chính thức Có thể Mở rộng cho Mạng Nơ-ron Sâu trong Các Hoạt động Chăm sóc Sức khỏe Quan trọng." (Tài trợ 3,2 triệu USD).
Sáng kiến Hệ thống Tự hành DARPA: "Học Tăng cường Đa Tác tử trong Môi trường Bị Hạn chế Truyền thông." (Tài trợ 2,8 triệu USD).
Thạc sĩ Khoa học Dữ liệu & Trí tuệ Nhân tạo — Sổ tay Học thuật
Trang 10 / 12


19. Tiêu chuẩn Tuyển sinh \& Quy trình Nộp hồ sơ
Chương trình MS-DSAI tìm kiếm các ứng viên có năng lực định lượng xuất sắc và nền tảng lập trình vững chắc. Việc tuyển sinh được thực hiện theo hình thức xét tuyển cuốn chiếu (rolling basis) qua các kỳ nhập học Mùa Thu và Mùa Xuân.

Điều kiện tiên quyết về Học thuật
Bằng Cử nhân ngành Khoa học Máy tính, Khoa học Dữ liệu, Toán học, Thống kê, Vật lý, hoặc Kỹ thuật.
Điểm GPA Đại học tối thiểu: 3.20 trên thang điểm 4.00.
Các môn học tiên quyết: Giải tích III, Đại số tuyến tính, Lý thuyết Xác suất, và Cấu trúc Dữ liệu bằng C++/Python.

Danh mục Hồ sơ Ứng tuyển
Bảng điểm Đại học chính thức.
3 Thư giới thiệu (Học thuật/Chuyên môn).
Bài luận Nguyện vọng (1.000 từ).
Bài thi GRE General (Không bắt buộc, khuyến nghị điểm định lượng > 165).
TOEFL iBT (Tối thiểu 100) hoặc IELTS (Tối thiểu 7.5) dành cho người không nói tiếng Anh bản ngữ.

20. Học phí, Lệ phí \& Trợ cấp Tài chính

Danh mục Chi phí | Chi phí mỗi Đơn vị | Tổng Chi phí Toàn chương trình (36 Tín chỉ)
Học phí Sau đại học | \$1,350 mỗi tín chỉ | \$48,600
Phí Cơ sở hạ tầng Siêu máy tính \& Phòng thí nghiệm | \$450 mỗi học kỳ | \$1,800
Phí Dịch vụ \& Sức khỏe Sinh viên Sau đại học | \$350 mỗi học kỳ | \$1,400
Ước tính Tổng Học phí \& Lệ phí | — | \$51,800 USD

Hỗ trợ Tài chính \& Học bổng Khuyến học Sau đại học
Trợ lý Nghiên cứu (RA)
Cung cấp 100\% học phí + trợ cấp hàng tháng 2.800 USD. Được trao dựa trên mức độ cạnh tranh và sự phù hợp về nghiên cứu với các dự án tài trợ của giảng viên. Yêu cầu 20 giờ/tuần làm việc tại phòng thí nghiệm.

Trợ giảng Giảng dạy (TA)
Cung cấp từ 50\% đến 100\% học phí + trợ cấp 2.400 USD/tháng. TA hỗ trợ giảng viên chấm điểm, dẫn dắt các buổi thực hành và tổ chức giờ tư vấn/hướng dẫn sinh viên.

Thạc sĩ Khoa học Dữ liệu \& Trí tuệ Nhân tạo — Sổ tay Học thuật
Trang 11 / 12

21. Định hướng Nghề nghiệp \& Việc làm Ngành Công nghiệp
Sinh viên tốt nghiệp chương trình MS-DSAI được tuyển dụng bởi các tập đoàn công nghệ hàng đầu thế giới, phòng thí nghiệm nghiên cứu, tổ chức tài chính và các công ty khởi nghiệp đổi mới sáng tạo có tốc độ tăng trưởng cao.

\$148k
MỨC LƯƠNG KHỞI ĐIỂM TRUNG BÌNH

98\%
CÓ VIỆC LÀM TRONG VÒNG 6 THÁNG

45+
ĐỐI TÁC TUYỂN DỤNG DOANH NGHIỆP

12\%
TIẾP TỤC HỌC CHƯƠNG TRÌNH TIẾN SĨ

Các Vị trí Tuyển dụng Hàng đầu \& Đối tác Doanh nghiệp Chính
Các Chức danh Công việc Mục tiêu Phổ biến
Kỹ sư Học máy Cao cấp
Kiến trúc sư Hệ thống AI
Nhà khoa học Dữ liệu Chính
Kỹ sư Cơ sở Hạ tầng MLOps
Chuyên gia Quản trị \& Pháp lý AI
Nhà nghiên cứu Thị giác Máy tính / NLP

Mạng lưới Đối tác Tuyển dụng
Google • Microsoft • Amazon Web Services • NVIDIA • Meta AI • OpenAI • Snowflake • Databricks • Goldman Sachs • JPMorgan Chase • Mayo Clinic Digital Health • Anthropic

22. Danh mục Kiểm tra Điều kiện Tốt nghiệp
Các Bước Hoàn thành Bằng cấp Cuối cùng:
1. Hoàn thành toàn bộ 36 tín chỉ bắt buộc với GPA tích lũy ≥ 3.00.
2. Nộp thành công và bảo vệ công khai thành công Luận văn / Đồ án Tốt nghiệp.
3. Xác nhận hoàn tất thủ tục chính thức từ Thư viện Trường và Cơ sở Siêu máy tính.
4. Nộp Đơn xin Tốt nghiệp với Phòng Đào tạo.

23. Thông tin Liên hệ của Học viện
Để biết thêm thông tin chi tiết về chương trình học thuật, yêu cầu tuyển sinh, hoặc hợp tác nghiên cứu, vui lòng liên hệ văn phòng chương trình:

Văn phòng Tuyển sinh Sau Đại học
Trường Máy tính, Học viện Công nghệ Apex
100 Innovation Boulevard, Tech District
Email: admissions-dsai@apex.edu

Phòng Giáo vụ Khoa
Điện thoại: +1 (555) 019-2831
Giờ làm việc: Thứ Hai – Thứ Sáu, 8:30 sáng – 5:00 chiều EST
Web: https://dsai.apex.edu/handbook

Thạc sĩ Khoa học Dữ liệu \& Trí tuệ Nhân tạo — Sổ tay Học thuật
Trang 12 / 12
\end{document}
